\section{Model-Feature Propensity Evaluation} \label{sec:propensity-eval}

To rigorously evaluate whether a model's \gls{fp} predictions reflected underlying phase-dependent seizure propensity, a set of qualification criteria was defined for each \textit{patient, feature,} and \textit{model}.
These criteria ensured that only meaningful model-feature relationships were interpreted in the context of cyclical seizure risk.

\subsection{Qualification Criteria} \label{subsec:qualification-criteria}
For a given patient, feature, and model, the following conditions must be satisfied for significance level $\alpha = 0.05$:

\begin{enumerate}
    \item \textbf{Seizure-Feature Phase Locking:} The patient's seizures must exhibit significant phase locking to the feature, as determined by the Rayleigh test with Benjamini-Hochberg \gls{fdr} correction ($q_{\mathrm{Rayleigh}} < \alpha$). This criterion ensures that the feature is even relevant to the patient's seizure timing.\\
    \textit{Again, due to limited seizure counts, seizures from both training and test sets are included, under the assumption that phase locking remains stable across the dataset.}
    \item \textbf{Model Predictive Performance:} The model must achieve a test set \gls{roc auc} $\geq 0.75$ and $q_{\mathrm{HanleyMcNeil}}<\alpha$.
    This empirically chosen threshold ensures that only models with acceptable predictive accuracy are considered.
    \item \textbf{\gls{fp}-Feature Phase Locking:} The model's \glspl{fp} must also be significantly phase locked to the feature ($q_{\mathrm{Rayleigh}} < \alpha$). This indicates that \glspl{fp} are not randomly distributed in phase.
    \item \textbf{Distributional Similarity:} If the above criteria are met, the phase distributions of \glspl{fp} and seizures are compared using the non-parametric permutation test (\cref{subsec:circular-comparison}). A non-significant result ($q_{\mathrm{Permutation}} \geq \alpha$) indicates that \glspl{fp} and seizures may arise from the same phase-dependent process (null hypothesis).
\end{enumerate}

\subsection{Interpretation of Outcomes} \label{subsec:interpretation-of-outcomes}

\paragraph{Ideal Case: Strong Evidence for Phase-Dependent Propensity}
\begin{itemize}
    \item Both seizures and \glspl{fp}: $q_{\mathrm{Rayleigh}} < \alpha$
    \item Permutation test: $q_{\mathrm{Permutation}} \geq \alpha$
\end{itemize}
\textit{Interpretation:} \Glspl{fp} occur at similar preferred phases as seizures, suggesting the model captures phase-dependent seizure propensity.

\paragraph{Contrary Cases}
\begin{itemize}
    \item For seizures, $q_{\mathrm{Rayleigh}} \geq \alpha$: seizures are not significantly phase-locked to the feature.
    Hence, any phase locking of \glspl{fp} cannot be interpreted as a meaningful proxy for propensity.
    \item For model \glspl{fp}, $q_{\mathrm{Rayleigh}} \geq \alpha$: \glspl{fp} are not significantly phase-locked to the feature.
    The model-feature combination cannot be used to evaluate propensity.
    \item Permutation test: $q_{\mathrm{Permutation}} < \alpha$: The phase preference of seizures and \glspl{fp} differs significantly.
    \Glspl{fp} cannot be interpreted as occurring during a proictal period.
\end{itemize}

As such, for each patient, model, and feature combination, it was evaluated whether \glspl{fp} indicated high propensity.